\section{Recall Forensics}
\label{sec:forensics}

Memory evaluations usually ask whether retrieval found a relevant item. A
different question matters after deployment: \emph{given a particular result
shown at a particular turn, what durable evidence explains why it appeared,
where it came from, and whether the evidence was later altered?} We call this
task \emph{recall forensics}. We have not identified a directly comparable
assistant-memory evaluation in the systems reviewed in \cref{sec:related};
this is a scoped observation, not a claim that no other system can implement
the task.

\plainbox{If the assistant said something surprising last Tuesday, the
database should show which accepted statement contributed to that result,
who supplied it, which score or connection selected it, and whether the
recorded trail still verifies. If the source was later deleted, the system
should prove the deletion rather than retain the words for convenience.}

\subsection{Task definition}
\label{sec:forensics-task}

The auditor receives the database, the evidence-format specification, and a
content-addressed source manifest. The reconstruction code does not call the
engine's ranking, graph, lifecycle, or audit helpers. For a logged result $X$
at recall turn $T$, it asks:

\begin{enumerate}
  \item[\textbf{F1}] \textbf{Attribution.} Recompute the text, trust, and
  recency components, Eq.~\ref{eq:record-score}, graph fusion when present,
  rank order, threshold decision, and result-limit decision from the bounded
  candidate ledger. A graph-nominated result must also carry a typed path
  whose objects were active according to graph events before $T$.
  \item[\textbf{F2}] \textbf{Provenance commitment.} Resolve record
  $\rightarrow$ admission $\rightarrow$ episode $\rightarrow$ origin and
  verify the chained content commitments. If a later deletion purged the
  source text, verify its prior digest and deletion transition without
  pretending the plaintext remains available.
  \item[\textbf{F3}] \textbf{Lifecycle consistency.} Replay admission,
  quarantine, explicit promotion, supersession, consolidation, and forgetting
  up to $T$. A returned record must be active at that point; the final row's
  status, source, trust tier, confidence, scope, and episode link must agree
  with the complete transition history.
  \item[\textbf{F4}] \textbf{Replay.} First, replay the recorded ranking
  decision for every turn. Separately report exact engine re-execution only
  for turns whose canonical inputs still exist unchanged. A turn before a
  later privacy deletion is deliberately ineligible for engine re-execution.
  \item[\textbf{F5}] \textbf{Integrity.} Test named mutation classes against
  copies of the database and report each class. This supports only the tested
  adversary model, not the universal statement that every possible forged
  field must be detected.
\end{enumerate}

Deterministic retrieval reduces the artifact set needed for F4; it does not
make embedding retrieval unreplayable in principle. A dense retriever can be
replayed if the exact model and tokenizer bytes, vector-index state,
preprocessing, numerical environment, dependencies, and code are preserved.
The narrower advantage here is that replay does not require preserving an
additional learned-model artifact.

\subsection{Evidence bindings}
\label{sec:forensics-binding}

Candidate summaries and paths live in a retrieval-event row rather than
directly in the audit chain. The evaluated post-v0.3.0 development snapshot
therefore stores a versioned evidence envelope and places its digest in the
chained \code{memory.recall} event:
\begin{equation}
D_{\mathrm{ret}}=\sha\!\left(\canon(v,i,s,u,q,q_{\mathrm{sha}},C,R,G)\right),
\label{eq:retrieval-binding}
\end{equation}
where $v$ is the format version, $i$ the retrieval identifier, $s$ and $u$
the subject and session, $q$ the retained query (empty by default),
$q_{\mathrm{sha}}$
its digest, $C$ the bounded candidate ledger, $R$ the ordered returned set,
and $G$ the graph and replay parameters. The envelope records one-based final
and lexical ranks, raw text-normalization inputs, trust tier, recency rank and
denominator, base score, graph rank and strength, and the typed path for each
returned graph nominee. The first 50 ranked candidates are retained, plus
every returned candidate when \code{limit} exceeds 50.

Admission events bind record content and the lifecycle metadata used by the
auditor: source type, initial trust and status, fact key, extraction
confidence, scope, and episode identifier. Promotion and later transitions
remain separate chained events. Episode ingestion already binds the original
message digest. These commitments detect direct changes to the selected
retrieval, record, and episode fields while preserving the rule that a
tombstone removes content.

The forensics experiment is not attributed to release commit
\code{12925d3} alone. That commit is only the snapshot's base. The evidence
file contains SHA-256 digests for 42 source files, including the evaluator and
the modified package, plus a digest over the manifest. Reproduction therefore
targets those bytes rather than an incorrect clean-commit claim.

\subsection{Protocol}
\label{sec:forensics-protocol}

The harness is stored at:
\begin{center}
\path{paper/scripts/run_recall_forensics_eval.py}
\end{center}
It creates nine initial scenario records, including an unreviewed web claim,
a second web-derived record that is explicitly promoted, a correction target,
a secret later deleted, and a two-hop relation, plus 1{,}000 filler records.
It executes three recall turns, then applies a correction and deletion,
executes eleven stable-state turns, and finally rebuilds the derived graph.
Across 14 turns it varies limits from 3 to 7, exercises a non-null threshold,
and enables graph recall in nine turns.

The standard-library reconstruction verifies all turns after the later
mutations and graph rebuild. Engine re-execution uses the logged
\code{limit}, \code{min\_score}, candidate cap, and graph-mode fields; it is
restricted to the eleven stable-state turns. The three earlier turns remain
decision-reconstructable but are excluded from engine re-execution because a
correction or deletion changed their canonical inputs. Fifteen selected
single-row mutation classes cover audit payload and event changes, retrieval
scores, returned sets, parameters, query and subject fields, record content
and lifecycle metadata, and episode text. Five trials are run per class.

\begin{table}[t]
\centering
\caption{Recall-forensics results for the content-addressed snapshot.}
\label{tab:forensics}
\small
\begin{tabularx}{\columnwidth}{@{}Yr@{}}
\toprule
Measure & Result \\
\midrule
Untampered chain, digest, lifecycle, turn failures & 0 \\
F1 ranking decisions recomputed & 14/14 \\
F1 returned-result attribution & 68/68 \\
\quad historical graph paths valid & 14/14 \\
\quad paths absent from current graph & 1 \\
F2 provenance commitments complete & 68/68 \\
\quad later-deleted source payloads unavailable & 1 \\
F3 lifecycle admissible at recall & 68/68 \\
\quad explicit promotion transition verified & yes \\
F4 stable returned sets re-executed exactly & 11/11 \\
\quad complete candidate ledgers exact & 11/11 \\
\quad historical turns excluded from re-execution & 3 \\
F5 selected single-row mutations detected & 75/75 \\
\quad tail truncation, chain only & 0/5 \\
\quad tail truncation, anchored checkpoint & 5/5 \\
Quarantined / post-deletion results returned & 0 / 0 \\
Forgotten content rows remaining & 0 \\
Auditor reconstruction time & 1.154\,s \\
\bottomrule
\end{tabularx}
\end{table}

\subsection{Results and boundaries}
\label{sec:forensics-results}

\Cref{tab:forensics} reports the pinned run. The auditor independently
recomputed score arithmetic, ranking, thresholds, and limits for all 14
turns, and reconstructed attribution, provenance commitments, and lifecycle
admissibility for all 68 returned-result pairs. All 14 typed graph paths were
valid against historical graph events. One old path no longer resolved in the
current derived graph after deletion and rebuilding, demonstrating why
historical evidence must not depend on current index rows. The explicitly
promoted record surfaced only after its promotion event.

All eleven turns whose canonical inputs remained stable reproduced both the
returned identifiers and complete bounded candidate ledgers. The three prior
turns were not counted as exact engine replays: retaining deleted plaintext
solely to make a benchmark replayable would violate the deletion contract.
For one returned-result pair, the later-deleted source is represented by its
ingestion digest, origin metadata, and deletion event; the source words are
correctly unavailable.

All 75 selected single-row mutations were detected. This is mutation coverage,
not a proof against arbitrary database rewriting. A coherent attacker who
rewrites the complete file and recomputes an unanchored chain remains outside
the file-only trust model. Tail truncation was correspondingly undetectable
from the chain alone (0/5) and detected against an anchored head (5/5).

Further limits are material. The candidate ledger is bounded; the auditor
recomputes selection within that sealed ledger but does not independently
regenerate every FTS candidate from historical corpus state. With query-text
retention disabled, it can confirm a proposed query against the digest but
cannot discover an unknown query. The workload is synthetic and single-host,
and the results establish neither fact truth nor general retrieval quality.
They establish what the measured evidence can and cannot reconstruct.
